\chapter{Conclusion}\label{sec:conclusion}
This paper investigated whether the Phillips curve relationship carries  forecasting value for Danish inflation over the period from December 2000 to March 2026, asking whether incorporating Danish unemployment and broader European macroeconomic variables could improve upon a univariate ARIMA  benchmark over a 24-month out-of-sample test period from April 2024 to March 2026.

The main finding is that the pre-surge ARIMA$(4,0,2)(0,0,1)_{12}$ outperforms every multivariate specification out-of-sample, achieving the lowest RMSE of $0.462$ and MAE of $0.351$. The Phillips Curve specification, despite conditioning on six lags of unemployment, attained an RMSE of $0.552$ and MAE of $0.490$, offering no incremental forecast content. Moreover, the model fails the Ljung--Box test, indicating remaining autocorrelation in the residuals. This result is consistent with the well-documented finding in the inflation forecasting literature that multivariate models often fail to outperform univariate benchmarks out-of-sample \parencite{stock_watson_2008}.

The GETS procedure provides a direct answer to the Phillips curve question. Starting from a general unrestricted model, the reduction eliminated every unemployment lag as statistically insignificant. Only contemporaneous EU27 inflation survived, confirming that the ERM\,II spillover channel is the dominant source of external information for Danish inflation. This result has a structural interpretation: under ERM\,II, monetary policy is subordinated to the exchange rate objective, which largely disables the domestic transmission mechanism through which unemployment would ordinarily dampen inflation. Danish inflation is instead anchored to the European price level, rendering domestic slack conditions a secondary determinant of inflation dynamics. However, while EU27 inflation is highly informative in-sample, it does not translate into improved out-of-sample forecasts: the GETS model achieves an RMSE of $0.546$ and MAE of $0.491$, marginally outperforming the naive ARIMA but falling well short of the pre-surge benchmark. This may partly reflect the absence of MA terms in the GETS specification. Because the GUM was specified with AR lags only, any moving-average dynamics in the error term are left unmodelled, which could weaken short-run forecast accuracy relative to the ARIMA benchmark. Beyond the specification, the EU27 relationship may also overfit to the unusual surge dynamics of the training period and not generalise to the subsequent disinflation.

Several limitations should be noted. The 24-month test period covers a single post-surge normalisation episode, and results may not generalise to other inflation regimes. The conditional forecasting approach for dynamic regression models assumes perfect foresight of future unemployment and EU27 inflation, overstating their practical forecasting value in real time; in particular, the contemporaneous EU27 channel presupposes that Eurostat figures are released prior to the Danish national release within the same month.

In sum, the evidence suggests that the Phillips curve has limited practical forecasting value for Danish inflation in the current sample. The domestic transmission mechanism is structurally constrained under ERM\,II, and past values of inflation itself remain the most reliable predictor - a finding that should be interpreted with the caveat that the test period covers a single post-surge normalisation episode. Future research could explore time-varying parameter models or nonlinear specifications that adapt more flexibly to structural change, or extend the analysis to other small open economies with similar monetary arrangements.